% Generated draft: requires mathematical and source review.
\documentclass[11pt]{article}
\usepackage[margin=25mm]{geometry}
\usepackage{fontspec}
\setmainfont{Latin Modern Roman}
\usepackage{amsmath,amssymb,mathtools,unicode-math}
\setmathfont{Latin Modern Math}
\usepackage{enumitem,needspace,xurl,hyperref}
\setlength{\parindent}{0pt}
\setlength{\parskip}{5pt}
\setlength{\emergencystretch}{3em}
\newcommand{\N}{\mathbb{N}}
\newcommand{\Z}{\mathbb{Z}}
\newcommand{\Q}{\mathbb{Q}}
\newcommand{\R}{\mathbb{R}}
\newcommand{\C}{\mathbb{C}}
\newcommand{\F}{\mathbb{F}}
\newcommand{\A}{\mathbb{A}}
\newcommand{\PP}{\mathbb{P}}
\newcommand{\E}{\mathbb{E}}
\newcommand{\eps}{\varepsilon}
\newcommand{\dd}{\,\mathrm{d}}
\begin{document}
\section*{UnsolvedMath ID: 30005386}
\textbf{Source problem code:} catalogue entry
\subsection{Definitions and mathematical statement}
Let \(K\subset S^3\) be a knot with tubular neighbourhood \(\nu K\), and write \(E(K)=S^3\setminus\operatorname{int}(\nu K)\) for its exterior. For the following formulation, use the conventional setting of tame, unoriented knots in the oriented sphere, with knots identified by ambient isotopy. These category and equivalence conventions are not explicitly specified in the supplied record and require review.

Choose a meridian \(\mu\) and preferred longitude \(\lambda\) on \(\partial E(K)\). A finite rational slope \(r=p/q\in\Q\), written with \(p\in\Z\), \(q\in\N\), and \(\gcd(|p|,q)=1\), is the unoriented isotopy class represented by \(p\mu+q\lambda\). Let \(S^3_K(r)\) denote the oriented closed three-manifold obtained by attaching a solid torus to \(E(K)\) so that its meridian is identified with this slope. Use the corresponding meridian--longitude convention for every comparison knot.

A finite rational slope \(r\) is \emph{characterising for \(K\)} if, for every knot \(K'\subset S^3\),
\[
S^3_K(r)\cong_{+}S^3_{K'}(r)
\quad\Longrightarrow\quad
K'\text{ is ambient isotopic to }K,
\]
where \(\cong_{+}\) denotes an orientation-preserving homeomorphism. It is \emph{non-characterising for \(K\)} if there exists a knot \(K'\) not ambient isotopic to \(K\) for which \(S^3_K(r)\cong_{+}S^3_{K'}(r)\). Set
\[
\mathcal C(K)=\{r\in\Q:r\text{ is characterising for }K\},
\qquad
\mathcal N(K)=\Q\setminus\mathcal C(K).
\]
\textbf{Classification problem (Question 3).} For every knot \(K\subset S^3\), determine \(\mathcal C(K)\) and \(\mathcal N(K)\); equivalently, give necessary and sufficient conditions on \(K\) and \(r\) for each alternative.

\textbf{Scope qualification.} The record does not explicitly define the slope domain or equivalence conventions. The formulation above uses finite rational slopes and orientation-preserving surgery comparisons; the latter convention is supported by the supplied example. A geometric slope can also be the meridian \(1/0\). Meridional filling returns \(S^3\) for every knot, so including that slope would be incompatible with the supplied assertion that every slope characterises the unknot. Thus exclusion of the meridian is an explicit interpretation requiring source review, rather than an additional verified source hypothesis. An unoriented homeomorphism convention could also change the classification.

\textbf{Supplied contextual results and their qualifications.} The following are reported background, not results proved or independently verified in this draft.
\begin{itemize}
\item The record attributes to Brakes a construction of non-characterising slopes by gluing two knot complements, citing report reference [3]. Its Example 4, citing [17], states that for all \(m,n,q\in\Z\setminus\{0\}\) with \(m\ne n\), there are distinct multiclasped doubles of double twist knots with orientation-preservingly homeomorphic \(1/q\)-surgeries. The extracted knot expressions are \texttt{W n (Tqm)} and \texttt{W m (Tqn)}. Their indices and the definitions of these knot families are not supplied clearly enough to reconstruct the notation reliably.
\item The record attributes to Baker and Motegi [2] examples with infinitely many non-characterising slopes. It identifies the simplest example as the hyperbolic knot extracted as \texttt{86}, for which no integer slope is characterising. The conventional interpretation of that extracted label is \(8_6\), but the typography requires confirmation.
\item The record reports that every slope is characterising for the unknot [7, 8], the trefoils, and the figure-eight knot [15], subject to the slope-domain qualification above. It also reports that all but finitely many non-integer slopes are characterising for torus knots [14, 11].
\item The extracted Lackenby assertion [9] says that for every knot \(K\) there is a constant \(C(K)\) such that every slope \(p/q\in[-1,1]\) is characterising, and says that Sorya refined this in [16]. As extracted, \(C(K)\) occurs in no condition on the slope. A missing inequality or other transcription defect is unresolved; no corrected bound is asserted here. The report describes the motivation as explicitly constructing such a lower bound in special cases.
\item The supplied Theorem 5, attributed to [17], concerns a satellite knot \(K=P(C)\) whose pattern complement has shortest-geodesic length at least \(0.136\), and reports that every slope satisfying
\[
\gcd(p,w)\ne1,\qquad |q|\ge36
\]
is characterising. The symbol \(w\), the precise pattern-complement construction, and the geometric structure used for geodesic length are not defined in the supplied excerpt. In particular, interpreting \(w\) as a winding number requires confirmation. This is retained as qualified background rather than a fully specified theorem.
\item The supplied Theorem 6, also attributed to [17], concerns a Whitehead double \(K=W(C)\). It assumes completeness of the SnapPea census of the first \(n\) orientable, two-cusped hyperbolic three-manifolds ordered by volume. Under that assumption, it reports that every slope with \(|p|\ne1\) and \(|q|\ge\phi(n)\) is characterising, with the supplied values
\[
\begin{array}{c|rrrrrrrrr}
n&4&6&8&10&16&18&24&49&289\\
\hline
\phi(n)&44&33&31&29&28&27&26&25&24
\end{array}
\]
The excerpt does not specify the Whitehead-double convention or define \(\phi\) beyond this table; no interpolation or extension is intended.
\end{itemize}
The supplied dated curation labels the problem partially solved, reports known characterising and non-characterising slopes in many families, and attributes to Sorya a sufficiently-large-denominator result for every knot. It states that a complete classification for each knot is not known. These are supplied literature assessments, not independently verified mathematical or current-status conclusions.
\subsection{Short English statement}
For each knot \(K\subset S^3\), classify the finite rational surgery slopes for which the oriented surgery manifold determines \(K\), and those for which a distinct knot has the same oriented surgery manifold. The slope domain, knot conventions, and several extracted background statements require review. The supplied partially solved status is dated curation, not an independently verified assessment.
\subsection{Sources}
\begin{itemize}
\item Supplied local UnsolvedMath registry record, permanent ID \(30005386\), title: \emph{Characterising and Non-Characterising Slopes of Knots}; exact source problem code: \texttt{catalogue entry}. No individual proposer is named in the supplied record.
\item \emph{Morphisms in Low Dimensions} (2023), Oberwolfach Reports, DOI \texttt{10.4171/owr/2023/4}, pp.\ 215--259, Question 3 and accompanying discussion. Supplied source URL and literature-source link labelled \emph{Oberwolfach Report 4/2023}: \url{https://doi.org/10.4171/owr/2023/4}. The registry says its clean statement was checked against the official report; this verification claim is supplied provenance and was not independently repeated here.
\item Sorya, \emph{Characterizing slopes for satellite knots}, supplied literature source: \url{https://doi.org/10.1016/j.aim.2024.109746}.
\item Sorya, \emph{Characterizing slopes for satellite knots}, supplied preprint link: \url{https://arxiv.org/abs/2307.00739}.
\item The excerpt attributes background constructions and results to Brakes [3], Baker and Motegi [2], Lackenby [9], and Sorya [16], and cites [7, 8], [15], [14, 11], and [17] for further statements. These are reference numbers in the source report; their full bibliography entries and their correspondence with the linked works are not supplied. They have not been reconstructed.
\item Literature assessment date: \texttt{2026-08-22}. The supplied dated curation labels the problem \emph{partially solved}, reports results in many families and a sufficiently-large-denominator result of Sorya for every knot, and states that a complete classification for each knot is not known. The background describes this as literature triage derived from the supplied OWR curation and requests independent verification of time-sensitive status and cited claims. This is supplied dated curation, not independently verified literature review.
\end{itemize}
\end{document}
